\chapter{Stage 8: The Cycle-Boundary Self-Improvement Loop}
\label{ch:cycle}

Everything in Part II happened without the model learning anything. A query was
triaged, routed, answered, verified, scored, and presented, and through all of it the
weights, the calibration, and the memory were only \emph{read}. This chapter is the
other clock: the slow, periodic loop that runs at \emph{cycle boundaries} and is the
only place the model actually changes. It is \caem{}'s answer to the third structural
flaw of \Cref{ch:problem}, staticness, and it is where the verified answers the
pipeline produced are folded back into the model so that next cycle's pipeline is
better than this one's.

\begin{intuition}[The system grades its own homework, then studies it, carefully]
The per-query pipeline does the system's homework all day and files the answers it
trusts into memory. At a cycle boundary the system stops taking new questions, opens
that file of verified answers, and \emph{studies them}: it fine-tunes itself on its own
best work. But studying can go wrong, so before keeping the lesson it checks that it
has not forgotten anything it used to know. If it has, it tears up the lesson and keeps
yesterday's self, while keeping the new file of answers. One clock answers; the other
clock learns, under a guard that refuses to let learning become forgetting.
\end{intuition}

\section{Two clocks, in full}
\label{sec:cycle-clocks}

\Cref{ch:bigpicture} introduced the two clocks; now we run the slow one in detail. The
fast clock is the per-query pipeline, ticking thousands of times. The slow clock is the
cycle: one tick processes a whole batch of queries through the fast clock and then runs
a single training-and-maintenance block. The realised run had six ticks of the slow
clock, cycles zero through five.

\begin{precise}[Why training is periodic, not per query]
Training on every query would be both ruinously expensive and dangerous: a gradient
step per answer would let one bad answer move the weights immediately, with no chance
to verify it first. Training never, and the system is static. \caem{} takes the middle
path: it \emph{accumulates} verified answers during the fast clock and trains on them in
a \emph{batch} at the cycle boundary, after they have all passed the gate. This is what
lets the loop train only on answers it has had a chance to trust, and what lets a single
guard check the whole update before committing it.
\end{precise}

\section{What one cycle does, in order}
\label{sec:cycle-order}

A cycle boundary is not one action but an ordered block, and the order is load-bearing.
\Cref{fig:cycle-master} lays out the full sequence; the sections that follow open each
part.

\begin{figure}[htbp]
\centering
\begin{tikzpicture}[
  font=\small\sffamily,
  inf/.style={rectangle, rounded corners=3pt, draw=cbIntuition!75!black, fill=cbIntuition!8,
    line width=0.9pt, align=center, inner sep=4pt, text width=10.5cm, font=\footnotesize, minimum height=0.7cm},
  trn/.style={rectangle, rounded corners=4pt, draw=cbExample!70!black, fill=cbExample!10,
    line width=1pt, align=center, inner sep=4pt, text width=10.5cm, font=\footnotesize, minimum height=0.7cm},
  fork/.style={diamond, aspect=2.2, draw=cbPitfall!75!black, fill=cbPitfall!10,
    line width=1pt, align=center, inner sep=1pt, font=\scriptsize, text width=2.6cm},
  bnd/.style={rectangle, rounded corners=4pt, draw=cbFormula!70!black, fill=cbFormula!8,
    line width=1pt, align=center, inner sep=4pt, text width=10.5cm, font=\footnotesize, minimum height=0.7cm},
  flow/.style={-{Stealth[length=2.2mm]}, line width=1pt, draw=cbInk!75},
  no/.style={-{Stealth[length=2.2mm]}, line width=1pt, draw=cbPitfall!75, dashed},
]
  \node[inf] (a) at (0,5.0) {\textbf{Score the calibration fold} under the current model};
  \node[trn] (b) at (0,4.1) {\textbf{Fine-tune (LoRA)} on collected verified episodes ($\ustored\ge0.70$)};
  \node[fork] (g) at (-3.0,2.9) {retention\\guard\\{\scriptsize worst probe $\ge 0.93$?}};
  \node[bnd] (c) at (0,1.7) {\textbf{Recalibrate}: refit temperature + composite, reload};
  \node[bnd] (d) at (0,0.8) {\textbf{Retro-verify} memory under the fresh calibration};
  \node[inf] (e) at (0,-0.1) {\textbf{Grow memory} with a fresh disjoint query chunk \;(always)};
  \node[inf] (f) at (0,-1.0) {\textbf{Evaluate} on held-out data, no storage \;(always)};
  \node[fork, draw=cbNumbers!75!black, fill=cbNumbers!12, text=cbNumbers!65!black] (h) at (0,-2.3)
    {early-stop?\\{\scriptsize 2 of 3 signals}};
  \node[bnd, text width=2.4cm, fill=cbInk!6, draw=cbInk!55] (revert) at (-7.2,2.9) {roll back to\\last good adapter};

  \draw[flow] (a) -- (b);
  \draw[flow] (b) -- (g);
  \draw[flow] (g) -- node[right,font=\scriptsize,text=cbExample!55!black]{pass (commit)} (c);
  \draw[no] (g) -- node[above,font=\scriptsize,text=cbPitfall!70!black]{fail (abort)} (revert);
  \draw[flow] (c) -- (d);
  \draw[flow] (d) -- (e);
  \draw[flow] (e) -- (f);
  \draw[flow] (f) -- (h);
  \draw[no] (revert) |- (e);
  \draw[flow] (h.south) -- ++(0,-0.5) node[below,font=\scriptsize,text=cbInk!70]{next cycle, or stop};
\end{tikzpicture}
\caption[One cycle, in order]{The cycle-boundary block. The model is fine-tuned on the
verified episodes, then the retention guard decides commit or abort. On a commit
(green) the calibration is refit, the memory re-verified, and the cycle proceeds; on an
abort (red) the adapter rolls back to the last good state and the commit-gated steps are
skipped. Growing the memory with fresh data and evaluating on held-out data run
regardless (blue). Finally the early-stop gate decides whether to begin another cycle. (Schematic.)}
\label{fig:cycle-master}
\end{figure}

\section{What it trains on: the verified episodes}
\label{sec:cycle-data}

The training data is the system's own memory, filtered. Only episodes whose stored trust
score $\ustored$ is at least $0.70$ are collected, a threshold deliberately \emph{higher}
than the $0.60$ storage cut of \Cref{ch:composite}: storage admits an answer to memory,
but training demands a stricter, cleaner subset. Of those, only episodes from the three
\emph{training} benchmarks are eligible; the two transfer benchmarks are held out so that
generalisation is measured on data the model never trained on.

\begin{precise}[The training target is the reasoning, not just the answer]
For each eligible episode the training target is its \emph{reasoning chain}, not merely
its final answer. The prompt is the question in the scaffolded chat format of
\Cref{ch:tiers}, and the target the model learns to produce is the verified
``Reasoning: \dots\ Answer: $Y$'' completion, with the prompt tokens masked out of the
loss so only the completion is learned. This is \emph{distillation}: the model is taught
to reproduce, from its own parameters, the grounded reasoning it previously needed
retrieval to produce. Chains that are empty, too short, or detected as repetitive loops
fall back to the bare answer or are filtered out, so the loop never trains on degenerate
text.
\end{precise}

\begin{figure}[htbp]
\centering
\resizebox{\textwidth}{!}{%
\begin{tikzpicture}[
  font=\small\sffamily,
  st/.style={rectangle, rounded corners=4pt, line width=1pt, align=center, inner sep=5pt,
    text width=2.6cm, minimum height=1.1cm, drop shadow={shadow xshift=0.3ex, shadow yshift=-0.3ex, fill=black!10}},
  flow/.style={-{Stealth[length=2.4mm]}, line width=1.1pt, draw=cbInk!75},
]
  \node[st, draw=cbExample!70!black, fill=cbExample!10] (mem) at (0,0)
    {\textbf{episodic memory}\\{\footnotesize all stored episodes}};
  \node[st, draw=cbIntuition!75!black, fill=cbIntuition!10] (col) at (3.7,0)
    {\textbf{collect}\\{\footnotesize $\ustored\ge0.70$, training benches, no loops}};
  \node[st, draw=cbFormula!70!black, fill=cbFormula!10] (rw) at (7.4,0)
    {\textbf{pool-reweight}\\{\footnotesize balance benches, cap $40\%$}};
  \node[st, draw=cbScenario!70!black, fill=cbScenario!10] (ft) at (11.1,0)
    {\textbf{LoRA fine-tune}\\{\footnotesize 3 epochs, CE on the chain}};
  \node[st, text width=2.0cm, draw=cbNumbers!75!black, fill=cbNumbers!12] (adp) at (14.3,0)
    {updated\\adapter};
  \draw[flow] (mem) -- (col);
  \draw[flow] (col) -- (rw);
  \draw[flow] (rw) -- (ft);
  \draw[flow] (ft) -- (adp);
\end{tikzpicture}}
\caption[The training-data flow]{Where the training set comes from. The memory is
filtered to high-trust, training-benchmark, loop-free episodes; the survivors are
re-weighted to balance the benchmarks; and the model is fine-tuned to reproduce their
verified reasoning chains. The output is an updated low-rank adapter. (Schematic.)}
\label{fig:cycle-data}
\end{figure}

\subsection*{Pool reweighting: keeping one benchmark from eating the pool}

Left alone, one benchmark can dominate the training pool. In an earlier version of the
system, fact-checking episodes came to make up the entire pool and the model collapsed
on the other benchmarks. Pool reweighting prevents this with four mechanisms working
together.

\begin{bythenumbers}[The four reweighting mechanisms]
\begin{itemize}[leftmargin=1.3em, itemsep=1pt]
  \item \textbf{Temperature-mixed proportions} ($T = 2.0$): the per-benchmark shares are
  softened, turning a tenfold gap in raw counts into roughly a threefold gap in training
  weight.
  \item \textbf{Upsample cap} ($3\times$): a scarce benchmark's samples are replicated at
  most three times, so balance never comes from runaway duplication.
  \item \textbf{A floor} ($50$ samples): every benchmark gets at least this many, even
  if smoothing would round it to nothing.
  \item \textbf{A hard ceiling} ($40\%$): no single benchmark may exceed forty percent of
  the final pool, with the excess water-filled to the under-represented ones.
\end{itemize}
A cold-start fallback seeds any benchmark with no verified episodes from a small set of
gold-labelled examples, so the loop can train even on the first cycles.
\end{bythenumbers}

\section{The fine-tune: a small adapter, a frozen model}
\label{sec:cycle-finetune}

The actual training touches only the low-rank adapter of \Cref{sec:prereq-lora}. The
backbone stays frozen; only the rank-thirty-two adapter on the seven projections, about
two percent of the parameters, receives gradients. The model is trained to reproduce the
collected chains by cross-entropy on the completion tokens.

\begin{lstlisting}[caption={The fine-tune, condensed from \texttt{caem/training/self\_improvement.py}.}, label={lst:finetune}]
loader = DataLoader(QADataset(train_pairs), batch_size=4, shuffle=True)   # micro-batch 4
optimizer = AdamW(adapter_parameters, lr=2e-4)                            # LoRA learning rate
for epoch in range(3):                                                    # epochs_per_cycle
    for step, batch in enumerate(loader):
        loss = cross_entropy(model(batch).logits, batch.labels)           # prompt masked -100
        # LoRA path: no L2 anchor -- the frozen base bounds drift on its own
        loss.backward()
        if (step + 1) % 4 == 0:                                           # grad_accum -> eff. batch 16
            clip_grad_norm(adapter_parameters, 1.0); optimizer.step(); optimizer.zero_grad()
\end{lstlisting}

\begin{pitfall}[Why there is no anchor term on the LoRA path]
A full fine-tune of every weight would risk drifting far from the starting model, so the
classic remedy adds an \emph{anchor} penalty to the loss, $(\lambda/2)\lVert\theta -
\theta_{\text{prev}}\rVert^2$, pulling the new weights toward the previous cycle's. On
the deployed low-rank path this term is \emph{switched off}, and deliberately so: the
backbone is frozen and the adapter is tiny, so the model simply \emph{cannot} drift far,
the bounded adapter is the anchor. The explicit penalty survives only in the
full-fine-tune fallback (with $\lambda = 0.01$), kept for ablation. Choosing a small
adapter over a penalised full update is the cleaner way to bound forgetting, and it is
why the loss in the listing is plain cross-entropy.
\end{pitfall}

\begin{figure}[htbp]
\centering
\begin{tikzpicture}[font=\small\sffamily]
\begin{axis}[
  name=poolax,
  width=7.4cm, height=5.0cm,
  ybar, bar width=11pt,
  xmin=0.4, xmax=5.6, xtick={1,2,3,4,5},
  ymin=0, ymax=5600,
  xlabel={cycle index}, ylabel={training-pool size (episodes)},
  tick label style={font=\scriptsize}, label style={font=\scriptsize},
  ymajorgrids, grid style={cbInk!8},
  title={(a) training-pool growth}, title style={font=\scriptsize\bfseries},
  nodes near coords, nodes near coords style={font=\tiny, /pgf/number format/1000 sep={}},
]
  \addplot[draw=cbExample!60!black, fill=cbExample!25] coordinates {
    (1,107) (2,1474) (3,2966) (5,5006) };
  \addplot[draw=cbPitfall!70!black, fill=cbPitfall!22] coordinates {
    (4,3624) };
\end{axis}
\begin{axis}[
  at={(poolax.east)}, anchor=west, xshift=1.1cm,
  width=7.4cm, height=5.0cm,
  xmin=0.6, xmax=5.4, xtick={1,2,3,4,5},
  ymin=0.04, ymax=0.30,
  xlabel={cycle index}, ylabel={final-epoch training cross-entropy},
  tick label style={font=\scriptsize}, label style={font=\scriptsize},
  ymajorgrids, grid style={cbInk!8},
  title={(b) training loss}, title style={font=\scriptsize\bfseries},
]
  \addplot[draw=cbScenario!70!black, line width=1.1pt, mark=*, mark size=1.8pt] coordinates {
    (1,0.258) (2,0.171) (3,0.109) (4,0.070) (5,0.091) };
  \addplot[only marks, mark=*, mark size=2.6pt, draw=cbPitfall!75!black, fill=cbPitfall!30] coordinates {
    (4,0.070) };
  \node[font=\tiny, text=cbPitfall!70!black, anchor=south] at (axis cs:4,0.070) {abort};
\end{axis}
\end{tikzpicture}
\caption[Per-cycle training trajectory]{The training trajectory of the actual deployed
run. Panel~(a): the training-pool size grows monotonically across the committed cycles
($107$, $1474$, $2966$, $5006$ at cycles one, two, three, and five), with the aborted
cycle four (red) drawing on $3624$ episodes. Panel~(b): the final-epoch training
cross-entropy falls steadily ($0.258$, $0.171$, $0.109$) and reaches its lowest value of
the entire run, $0.070$, at cycle four (red), before rising to $0.091$ at cycle five.
Cycle four thus reached the lowest training loss of the run yet breached the retention
floor, the textbook signature of an over-fit update that the guard correctly caught and
rolled back.}
\label{fig:cycle-traj}
\end{figure}

\section{The retention guard: commit or abort}
\label{sec:cycle-guard}

A fine-tune can improve the target benchmarks while quietly eroding unrelated abilities,
which is the catastrophic forgetting of \Cref{ch:problem}. The retention guard is the
check that catches this before the update is kept.

\begin{precise}[Pristine probes, worst-case ratio, and the floor]
At cycle zero, before any training, the system measures its accuracy on three
\emph{retention probes}, held-out tasks chosen to span different abilities: a
four-option general-knowledge test, an open-text factoid test, and a five-option
commonsense test, two hundred questions each. These cycle-zero scores are the
\emph{pristine baseline}, anchored once and reused forever as the denominator (it is even
persisted to disk so a resumed run does not re-anchor against an already-trained model).
After each fine-tune the three probes are re-run, and each post-training score is divided
by its pristine value to get a retention ratio. If the \emph{worst} of the three ratios
falls below the floor $\rhomin = 0.93$, the cycle \emph{aborts}. The rule is a logical
AND over the probes: all three must hold above the floor for the cycle to commit, so a
collapse on any single ability vetoes the update.
\end{precise}

\begin{figure}[htbp]
\centering
\begin{tikzpicture}
\begin{axis}[
  width=10.5cm, height=5.6cm,
  ybar, bar width=20pt,
  symbolic x coords={general-knowledge, open-text factoid, commonsense},
  xtick=data, x tick label style={font=\scriptsize},
  ymin=0.7, ymax=1.08, ylabel={retention ratio (post / pristine)},
  ytick={0.7,0.8,0.9,0.93,1.0}, font=\small\sffamily, tick label style={font=\scriptsize},
  ymajorgrids, grid style={cbInk!8},
  nodes near coords, nodes near coords style={font=\scriptsize, fill=white, inner sep=1.5pt, /pgf/number format/fixed, /pgf/number format/precision=3},
]
  \draw[cbPitfall!75!black, dashed, line width=1.2pt] (axis cs:general-knowledge,0.93) -- (axis cs:commonsense,0.93);
  \node[font=\scriptsize, text=cbPitfall!70!black, anchor=east] at (axis cs:commonsense,0.945) {floor $\rho_{\min}=0.93$};
  \addplot[draw=cbExample!60!black, fill=cbExample!25] coordinates {
    (general-knowledge,1.018) (open-text factoid,0.886) (commonsense,0.982) };
\end{axis}
\end{tikzpicture}
\caption[The retention guard firing at cycle four]{The three retention probes measured
at the cycle-four boundary, read from the run log. General knowledge actually improved
slightly (ratio $1.018$) and commonsense held at $0.982$, both above the floor, but the
open-text factoid probe fell to $0.886$, below the $0.93$ floor. Because the worst probe
breaches the floor, the cycle aborts: the freshly trained adapter is discarded and the
previous cycle's adapter is restored. The memory accumulated this cycle is kept.}
\label{fig:cycle-guard}
\end{figure}

\begin{precise}[Asymmetric rollback: weights revert, memory persists]
On an abort the adapter is rolled back to the previous cycle's snapshot, and the
commit-gated steps below are skipped. But the \emph{episodic memory is not rolled back}:
every episode admitted during the aborted cycle survives. This asymmetry is the core
safety property of the whole architecture. Memory growth is monotone and high-precision
by construction, since each episode passed the gate, so it is always safe to keep; a
weight update is a global change with no such guarantee, so it is the thing thrown away.
The realised run hit this fork for real at cycle four: the open-text factoid probe fell to
$0.886$, below the floor, the adapter reverted to its cycle-three state, and the memory
carried forward untouched. The next cycle then re-attempted training from the last good
adapter against fresh data.
\end{precise}

\section{The rest of the boundary, and what is gated on committing}
\label{sec:cycle-boundary-ops}

If the cycle commits, four more things happen, in order, and each depends on the model
having actually improved.

\begin{description}[leftmargin=0pt, labelsep=0.5em, style=nextline]
  \item[Recalibrate.] The calibration fold is re-scored under the new model, then the
  temperature of \Cref{ch:upre} and the composite of \Cref{ch:composite} are re-fitted
  and reloaded. This is the per-cycle refit of \Cref{sec:composite-training}: it keeps
  $\ustored$ honest as the model moves underneath it.
  \item[Re-verify the memory.] Every stored episode is re-scored through the freshly
  recalibrated verifier. An episode whose trust has risen above the storage threshold is
  promoted, and one that has fallen below a floor is pruned. In the deployed run a typical
  committed cycle re-scored several thousand stored episodes and pruned only a small
  single-digit-percentage tail; at the third cycle, for instance, all $5358$ stored
  episodes were re-scored, $352$ were pruned, and $5006$ survived. This is how the deferred
  fact-checking claim of $\ustored = 0.508$ from \Cref{ch:bigpicture} is finally promoted to
  memory: the better model scores it higher, and it crosses the line. Re-verification runs
  \emph{after} recalibration, deliberately, so it scores through the new curves rather
  than stale ones. The aborted cycle skipped re-verification entirely, because re-scoring
  through unchanged weights would be a no-op.
  \item[Reconsider deferrals.] The deferred buffer of borderline answers is re-examined
  and each entry is promoted, expired, or kept.
  \item[Consolidate.] Near-duplicate episodes are merged into a single representative.
\end{description}

\begin{pitfall}[An abort skips the commit-gated half, but not the whole cycle]
All four of the above, plus the re-verification, are gated on the cycle committing; an
abort skips them, because re-scoring and re-verifying through an unchanged model would add
nothing. But the cycle still does its \emph{unconditional} work even on an abort: it grows
the memory with a fresh, content-hash-disjoint chunk of new questions (with storage on),
runs the held-out evaluation (with storage off), checkpoints, and consults the early-stop
gate. An aborted cycle loses its weight update, not its progress.
\end{pitfall}

\section{When the loop stops}
\label{sec:cycle-stop}

The loop is configured for ten cycles, but it does not always use them. After a burn-in,
an early-stop gate watches three signals: whether the joint quality score has stopped
rising, whether the storage rate has saturated, and whether general-knowledge retention is
holding. When two of the three agree that the system has reached equilibrium, the loop
stops early. In the deployed run the two signals that fired were the saturated storage
rate and the holding general-knowledge retention, while the joint-quality saturation
signal had not yet tripped; two of the three was enough to stop at the first eligible
cycle after the burn-in. The realised run stopped this way at cycle five, having committed
cycles one, two, three, and five and aborted cycle four, a commit fraction of four in
five. The full realised trajectory is the subject of \Cref{ch:trajectory}.

\section{Why this closes the static flaw}
\label{sec:cycle-why}

\begin{intuition}[Verification that compounds instead of evaporating]
\Cref{ch:problem} argued that the deepest waste in a static model is that verification
does not compound: effort spent confirming an answer today changes nothing tomorrow. The
cycle loop is the mechanism that makes it compound. Every verified answer is filed in
memory; every cycle the model studies that file and folds the knowledge into its
parameters; and the retention guard ensures each such step is a net gain rather than a
trade. The receipt that this is distillation and not mere caching is that, after the loop
runs, the model answers from its own parameters, on the zero-shot tier, questions it
previously could only answer with retrieval, and it does so at a higher accuracy than the
retrieval tier reached at cold-start: in the deployed run the learned zero-shot tier
reaches an exact match of about $0.57$ on its assigned queries by the third cycle, above
the roughly $0.48$ the retrieval tier managed on its assigned queries at the cold start.
The loop turns one-shot verification into a permanent capability.
\end{intuition}

\begin{bythenumbers}[Cycle-loop constants]
\begin{itemize}[leftmargin=1.3em, itemsep=1pt]
  \item Training set: episodes with $\ustored \ge 0.70$, training benchmarks only, loops filtered.
  \item Pool reweighting: temperature $2.0$, upsample cap $3\times$, floor $50$, ceiling $40\%$.
  \item Fine-tune: LoRA rank $32$, $\alpha = 64$, $3$ epochs, micro-batch $4$, effective batch $16$, learning rate $2\times10^{-4}$, AdamW, no anchor term.
  \item Retention guard: three probes ($200$ each), worst-ratio floor $\rhomin = 0.93$; abort restores the adapter, memory persists.
  \item Boundary: recalibrate, re-verify, reconsider, consolidate (all commit-gated); grow and evaluate (unconditional).
  \item Horizon: up to $10$ cycles; early-stop on two of three equilibrium signals (realised stop at cycle five).
\end{itemize}
\end{bythenumbers}

\begin{defence}[If an examiner asks: ``what shows the loop helps rather than just drifts?'']
Three guards make the answer empirical, not hopeful. The training set is gated at a strict
trust threshold, so the loop studies only high-quality answers. The retention guard
measures held-out ability after every update and rolls back any update that regresses it,
so a drifting cycle is caught and undone rather than committed. And the memory is curated
against the improving model each cycle, so its quality is maintained rather than left to
decay. The realised commit fraction of four in five, with one genuine abort, is the direct
evidence that the guard is active and that the committed cycles are the ones that earned
their keep.
\end{defence}

The per-query pipeline and the cycle loop are now both fully open. \Cref{ch:theory} steps
back to the theory: the data-purity property that lets an imperfect verifier yield a
clean memory, and the monotonicity and convergence results that the loop's behaviour is
meant to exhibit.
